La proliferación de información falsa en medios digitales se ha convertido en uno de los desafíos más apremiantes de la sociedad contemporánea. En el contexto de América Latina, y particularmente en Colombia, la desinformación sobre el sistema financiero genera consecuencias económicas directas: rumores sobre quiebras bancarias provocan retiros masivos, noticias falsas sobre congelamiento de cuentas generan pánico financiero, y fraudes disfrazados de comunicados oficiales causan pérdidas patrimoniales a ciudadanos vulnerables.

A pesar de la gravedad del problema, las herramientas actuales de detección de desinformación presentan dos limitaciones fundamentales. Primero, la gran mayoría de los recursos y modelos han sido desarrollados para el idioma inglés: el 81,5\% de los datasets existentes para detección de noticias falsas están en este idioma~\cite{Oshikawa2020}. Segundo, los enfoques basados exclusivamente en modelos de lenguaje de gran escala (LLMs) alcanzan precisiones de apenas 64--71\% en tareas de verificación de hechos~\cite{Hoes2023}, lo que resulta insuficiente para dominios donde los errores tienen consecuencias económicas reales.

El presente anteproyecto propone un sistema híbrido que aborda ambas limitaciones. El sistema combina la capacidad de clasificación de LLMs con un agente inteligente que consulta fuentes institucionales del ecosistema financiero colombiano (Superintendencia Financiera, Banco de la República, DANE, ColombiaCheck) cuando el modelo tiene baja confianza en su predicción. Esta arquitectura de activación condicional permite optimizar recursos computacionales al mismo tiempo que mejora la precisión y, fundamentalmente, la explicabilidad del sistema: cada veredicto incluye las fuentes consultadas y el razonamiento seguido por el agente.

El proyecto contempla tres contribuciones principales: (1) la construcción del primer benchmark de evaluación de desinformación financiera colombiana, (2) una arquitectura con umbral de confianza como detonador de un agente secundario de verificación, y (3) herramientas de consulta específicas para el ecosistema regulatorio financiero colombiano que no existen en la literatura actual.

El documento se organiza de la siguiente manera: el Capítulo 2 presenta el estado del arte en detección de desinformación y sistemas multi-agente; el Capítulo 3 formula el problema de investigación; el Capítulo 4 justifica la pertinencia del proyecto; el Capítulo 5 define los objetivos; el Capítulo 6 detalla la metodología; el Capítulo 7 presenta el cronograma; el Capítulo 8 estima el presupuesto; y el Capítulo 9 describe los resultados esperados.
